% ---------------------------------------------------------------------
%  Données et licences
% ---------------------------------------------------------------------
\chapter{Données et licences}\label{ann:donnees}

\section{Les données Olist}

Le jeu \emph{Brazilian E-Commerce Public Dataset by Olist} \citep{olist2018} décrit environ
100\,000 commandes passées entre 2016 et 2018 sur plusieurs places de marché brésiliennes. Les
données sont réelles et anonymisées : les noms des entreprises et des clients ont été remplacés.

\begin{center}
\small
\renewcommand{\arraystretch}{1.25}
\begin{tabularx}{\textwidth}{@{}l X l@{}}
\toprule
\textbf{Fichier} & \textbf{Contenu} & \textbf{Utilisé dans} \\
\midrule
\cmd{orders}, \cmd{order\_items} & commandes et lignes de commande & TP 1 à 4 \\
\cmd{order\_payments} & paiements & TP 1 \\
\cmd{customers}, \cmd{sellers} & clients et vendeurs, avec leur ville et leur état & TP 1 à 3 \\
\cmd{products}, \cmd{product\_category\_name\_translation} & produits et catégories, avec leur nom anglais & TP 1 à 3 \\
\cmd{order\_reviews} & avis des clients, en portugais & TP 2 et 3 \\
\bottomrule
\end{tabularx}
\end{center}

Les tables de vente sont chargées dans une base PostgreSQL, qui joue le rôle de l'application de
vente. Les autres sont fournies en fichiers CSV.

Le jeu est diffusé sous licence \textbf{Creative Commons Attribution -- Pas d'utilisation
commerciale -- Partage dans les mêmes conditions 4.0} (CC BY-NC-SA 4.0).

\section{Les données reconstituées}

Deux sources des travaux pratiques ont été reconstituées à partir des commandes Olist, pour
illustrer des situations réelles que le jeu d'origine ne contient pas :
\begin{itemize}[leftmargin=*,itemsep=2pt]
  \item \textbf{les messages de suivi des colis} (TP 2) : un message JSON par colis, avec ses étapes
        et, à partir de 2018, ses incidents. Les dates suivent celles des commandes. Quelques défauts
        sont ajoutés volontairement : messages tronqués, envoyés deux fois, changement de format ;
  \item \textbf{le flux d'événements} (TP 4) : les mêmes messages, découpés en événements et rejoués
        en temps accéléré, avec 5~\% d'événements en retard.
\end{itemize}
Ces données sont produites par les scripts du dépôt d'accompagnement, à partir de tirages
aléatoires à graine fixe : elles sont identiques d'une exécution à l'autre. Elles suivent la
licence du jeu Olist dont elles dérivent.

\section{Les photos de produits}

Les photos et les fiches du TP 3 sont un sous-ensemble du jeu \emph{Amazon Berkeley Objects}
\citep{collins2022}, publié par Amazon et l'université de Californie à Berkeley sous licence
\textbf{Creative Commons Attribution -- Pas d'utilisation commerciale 4.0} (CC BY-NC 4.0). Environ
3\,000 photos de six types de produits ont été retenues, réduites à 256 pixels au plus, et réparties
entre douze vendeurs fictifs. Des défauts ont été ajoutés : fichiers coupés, fichier qui n'est pas
une image, vignettes minuscules, photos envoyées deux fois ou recadrées.

\section{Le modèle de vision}

Le TP 3 utilise le réseau MobileNetV3 \citep{howard2019}, pré-entraîné sur le jeu ImageNet, converti
au format ONNX. Il sert uniquement à calculer des vecteurs, sans réentraînement.

\section{Conditions d'utilisation}

Ces licences autorisent la reproduction et l'adaptation des données à des fins non commerciales, à
condition de citer leur source. Cet ouvrage, à caractère et à usage pédagogiques, respecte ces
conditions. Toute réutilisation commerciale des données demanderait l'accord de leurs détenteurs.
